\documentclass[11pt]{article}

% basic packages
\usepackage{graphicx}
\usepackage[letterpaper, 
hmargin={0.7in,0.7in}, 
vmargin={0.7in,0.7in}]{geometry}

\usepackage{soul}
\usepackage{paralist,revnum}
\usepackage{microtype}
\usepackage{booktabs}
\usepackage{ragged2e}

\usepackage[scaled=0.85]{Univers}
%\usepackage{WarnockPro}         % also try ChaparralPro
\usepackage{ChaparralPro}
\usepackage{marvosym}

% tweaks
\setlength{\overfullrule}{5pt}
\setlength{\parindent}{0em}
\newcommand{\til}{\char158}
% \newcommand{\mybullet}{\ornament{19}} % 6, 8 are flashier % for warnock
\newcommand{\mybullet}{\scalebox{0.5}{\ornament{20}}}
\newcommand{\blankline}{\quad\pagebreak[2]}

\newcommand{\mysection}[1]{%
  \currentpdfbookmark{#1}{#1}
  \noindent{\Large \color{webbrown} \textsc{#1}}\\[-2ex]
  {\color{halfgray} \hrule}
  \vspace*{1ex}}

% An itemize-style list with lots of space between items
\newenvironment{outerlist}[1][\textbullet]%
        {\begin{enumerate}[#1]}{\end{enumerate}%
         \vspace{-.6\baselineskip}}

\newenvironment{innerlist}[1][\enskip\mybullet]%
{\begin{compactitem}[#1]}{\end{compactitem}}

% links and color
\usepackage{hyperref}
\usepackage[dvipsnames]{xcolor}

\definecolor{halfgray}{gray}{0.55}
\definecolor{medgray}{gray}{0.3}
\definecolor{webgreen}{rgb}{0.10,0.29,0.19}
\definecolor{webbrown}{rgb}{.5,.2,0}
\definecolor{darkblue}{rgb}{0.0,0.0,0.4}

\hypersetup{colorlinks,           
  breaklinks=true,
  linkcolor=darkblue,
  urlcolor=webgreen,
  anchorcolor=darkblue,
  citecolor=darkblue,
  pdfauthor={Michael McCourt},
  pdftitle={curriculum vitae},
  pdfcreator={pdflatex}}

\newcommand{\daterange}[2]{%
\begin{minipage}[t]{1in}
{\raggedleft 
\textsc{#1}--- \\ 
\textsc{#2}\phantom{---}

}
\end{minipage}
}

\newcommand{\datealt}[1]{%
\begin{minipage}[t]{1in}
{\raggedleft 
\textsc{#1}\phantom{---}

}
\end{minipage}
}

\newcommand\email{\href{mailto:mike@sturdystatistics.com}%
  {mike@sturdystatistics.com}}
\newcommand\web{\href{https://mikemccourt.com}%
  {mikemccourt.com}}
\newcommand\github{\href{https://github.com/mkmcc}%
  {github.com/mkmcc}}
\newcommand\arxiv{\href{https://ui.adsabs.harvard.edu/public-libraries/mwmNbnbJTXqQZahtQBn9BQ}%
  {arxiv.org/a/mccourt\_m\_1}}



% institutions
\newcommand\stanford{\href{http://www.stanford.edu/}%
  {Stanford University}}
\newcommand\bstanford{\href{http://www.stanford.edu/}%
  {\textit{Stanford University}}}

\newcommand\berkeley{\href{http://www.berkeley.edu/}%
  {UC Berkeley}}
\newcommand\bberkeley{\href{http://www.berkeley.edu/}%
  {\textit{UC Berkeley}}}

\newcommand\glam{\href{http://0-www.stanford.edu.ilsprod.lib.neu.edu:80/group/glam/}%
  {Geballe Laboratory for Advanced Materials}}

\newcommand\kipac{\href{http://www-group.slac.stanford.edu/KIPAC/}%
  {Kavli Institute for Particle Astrophysics and Cosmology}  (KIPAC)}

\newcommand\slac{\href{http://slac.stanford.edu}%
  {Stanford Linear Accelerator Center} (SLAC)}

\newcommand\tac{\href{http://astro.berkeley.edu/tac/}%
  {Theoretical Astrophysics Ceneter}}



% actual document
%%%%%%%%%%%%%%%%%%%%%%%%%%%%%%%%%%%%%%%%%%%%%%%%%%%%%%%%%%%%%%%%%%%%%%
\begin{document}
\raggedbottom
\pagestyle{empty}

% title and contact
%%%%%%%%%%%%%%%%%%%%%%%%%%%%%%%%%%%%%%%%%%%%%%%%%%%%%%%%%%%%
\newcommand{\contact}{%
%\Telefon \, 310-923-2656 \quad
\Telefon \, 123-456-7890 \quad
\Letter \, \email \quad
\ComputerMouse \, \web \\
\raisebox{-3pt}{\includegraphics[height=11pt]{blacktocat.pdf}} \github
\,\quad
\raisebox{-3pt}{\includegraphics[height=11pt]{arxiv.pdf}} \arxiv}

\newlength{\titlewidth}
%\settowidth{\titlewidth}{\contact}
\setlength{\titlewidth}{\linewidth}

\hfill
\begin{minipage}{\titlewidth}
  \centering
  {\huge 
    \color{webbrown} 
    \textls*[300]{\textsc{Michael{\,\,}K.{\,\,}McCourt}}} \\[2ex]
  %% {\color{medgray} 
  %%   Harvard-Smithsonian Center for Astrophysics,
  %%   60 Garden Street, 
  %%   Cambridge, MA 02138} \\
  {\color{medgray} %
    \contact}
\end{minipage}%
\hfill

\vspace*{\baselineskip}

%% \begin{center}
%%   \textit{Astrophysics PhD with expertise in programming and
%%     predictive modeling of complex systems.\\  Seeking technical or
%%     scientific role in an environmentally-focused company.}
%% \end{center}

%% \vspace*{\baselineskip}



%% Skills
%%%%%%%%%%%%%%%%%%%%%%%%%%%%%%%%%%%%%%%%%%%%%%%%%%%%%%%%%%%%
%
\mysection{Skills and Technical Proficiencies}
%
\begin{outerlist}
\item[] \textbf{Modeling}
\begin{innerlist}
  \item I specialize in building simple but predictive models for
    complex systems, identifying and highlighting the most salient
    features.  \textit{My goal is to make difficult problems easy.}
  \item Expertise in translating ideas rapidly into equations (or
    code) and vice-versa.
\end{innerlist}

\item[] \textbf{Programming}
\begin{innerlist}
  \item Proficient in C, Ruby, and Python.  Also proficient with
    Mathematica, Emacs Lisp, and shell scripts.
  \item Proficient in high-performance computing (utilizing some of
    the nation's fastest supercomputers), and in statistical data
    fitting using machine learning algorithms, genetic algorithms, and
    Markov Chain Monte Carlo simulations.
  \item Created popular open-source tools for simplifying daily
    scientific tasks, including:
    \begin{innerlist}
      \item[-]
        ``\href{http://mkmcc.github.io/software/gnuplot-mode.html}{gnuplot-mode},''
        a package for editing plots, downloaded by over 30,000
        scientists to date, and
      \item[-]
        ``\href{http://mkmcc.github.io/software/bibslurp.html}{bibslurp},''
        a tool for automatically making bibliographies from the
        standard NASA database, adopted by 2,000+ scientists
    \end{innerlist}
\end{innerlist}

\item[] \textbf{Communication and Collaboration}
\begin{innerlist}
  \item Invited speaker to 29 scientific conferences and seminars
  \item Colleagues often send me their proposals for comments, and one
    of my proposals for supercomputer time was selected by the
    granting agency as their published example of a well-written
    application
  \item Selected as a referee (expert reviewer) for five academic
    journals, including \textit{Nature Letters}
  \item Initiated several new collaborations among scientists
\end{innerlist}

\item[] \textbf{Management Experience}
\begin{innerlist}
  \item Raised over \$2 million 
    %(\$1.3 million as the sole or primary author)
    in funding and supercomputer time to run a small research group
  \item Conducted an international job search and hired a postdoctoral
    research scholar
  \item Mentored two undergraduate students and three graduate
    students on multi-year projects
  \item Organized and taught a fluid dynamics ``bootcamp'' at UCSB
    after recognizing it was needed
\end{innerlist}

\end{outerlist}

\blankline

%% Scientific Career
%%%%%%%%%%%%%%%%%%%%%%%%%%%%%%%%%%%%%%%%%%%%%%%%%%%%%%%%%%%%
%
\mysection{Experience and Education}
%

\begin{innerlist}
  \item Authored 25 peer-reviewed scientific papers cited 2,200
    times (research index quotient in the 95\textsuperscript{th}
    percentile for astrophysics at the time of my departure in 12/2017)
  \item Conducted self-directed research on 10 different topics over
    11 years
  \item Researching deep learning image recognition algorithms as a
    freelance consultant for a cybersecurity firm
\end{innerlist}

\blankline

\textbf{KITP Fellow}, UC Santa Barbara\hfill\textit{(awarded for 2019-2020)}\\
\textbf{Hubble Fellow}, NASA\hfill\textsc{Sept 2016}---present\\
\textbf{ITC Fellow}, Harvard University\hfill\textsc{Sept 2014---July 2015}
%
\begin{innerlist}
\item {\RaggedRight Recipient of prestigious independent research fellowships.}
\end{innerlist}
%

\vspace*{\baselineskip}

\textbf{Graduate Student Researcher}, UC Berkeley.\hfill\textsc{{Aug 2008---Aug 2014}}
%
\begin{innerlist}
\item {\RaggedRight Recipient of both department awards for
  outstanding graduate student research.}
\end{innerlist}
%

\vspace*{\baselineskip}

\textbf{PhD, Astrophysics}, UC Berkeley\hfill\textsc{{May 2014}}
%

\textbf{MA, Astrophysics}, UC Berkeley\hfill\textsc{{May 2010}}
%

\textbf{BS, Physics}, Stanford University (concentration in
theoretical physics)\hfill\textsc{{June 2008}}
%
%


% Postamble
%%%%%%%%%%%%%%%%%%%%%%%%%%%%%%%%%%%%%%%%%%%%%%%%%%%%%%%%%%%%
\enlargethispage{4\baselineskip}
\vfill
\begin{minipage}[h!]{\linewidth}
  \hspace*{\fill}
  \begin{center}
    {\scriptsize  Last updated: \today \\*[-0.5\baselineskip]
      current version: 
      \href{https://github.com/mkmcc/resume/raw/master/resume.pdf}%
      {github.com/mkmcc/resume}}
  \end{center}
  \hspace*{\fill}
\end{minipage}



%% \clearpage

%% % Research Achievements
%% %%%%%%%%%%%%%%%%%%%%%%%%%%%%%%%%%%%%%%%%%%%%%%%%%%%%%%%%%%%%
%% \newgeometry{margin=1in}
%% \mysection{Key Research Accomplishments}
%% %
%% \vspace*{\baselineskip}
%% %
%% Most of my academic research has focused on galaxy clusters, which are
%% the largest and most recent structures to form in the universe.  For
%% more information and visuals, visit my website:
%% \href{http://mkmcc.github.io/research/index.html}{mkmcc.github.io/research}

%% \begin{outerlist}
%%   \item Showed that clouds of cold gas can condense like rain in the
%%     hot atmospheres of galaxy clusters; this process fuels the growth
%%     of galaxies in clusters and determines the maximum mass of
%%     galaxies in the universe.  My research overturned a paradigm about
%%     thermal instability which had been in place for more than 20
%%     years. (read more \href{http://mkmcc.github.io/research/#sec:ti}{here}.)
%%   \item Discovered that gas in galaxy clusters may be vigorously
%%     ``boiling,'' contrary to what had been found in previous studies.
%%     This result influences how x-ray telescopes infer the masses of clusters,
%%     which are then used to determine quantities such as the age of the
%%     universe.  (read more
%%     \href{http://mkmcc.github.io/research/#sec:mti}{here}.)
%%   \item Showed that the unexplained range of temperatures inside
%%     galaxy clusters is in fact determined by variations in their
%%     formation histories. (read more \href{http://mkmcc.github.io/research/#sec:outer-parts}{here}.)
%%   \item Showed that clouds of cold gas in the center of the Milky Way
%%     Galaxy may survive longer than had previously been thought.
%%     Outlined a method for using these clouds to study the unknown
%%     physics of gas falling into the super-massive black hole in our
%%     galactic center. (read more
%%     \href{http://mkmcc.github.io/research/#sec:g2}{here}.)
%%   \item Showed that the bizarre and unexpected orbits of Pluto-like
%%     objects in our Solar System may indicate the presence of a second
%%     Kuiper belt --- this is a prediction which should be tested soon!
%%     (read more
%%     \href{http://mkmcc.github.io/research/index.html#sec:plutos}{here}).
%%   \item Showed that the cold gas recently discovered in the outskirts
%%     of galaxies is in fact a ``fog'' of tiny, dispersed chunks of gas
%%     rather than a contiguous, solid object.  Discovered a new
%%     hydrodynamic instability which explains why this should be the
%%     case, simultaneously explaining a number of observational puzzles
%%     in studies of galaxies.
%% \end{outerlist}

%% \restoregeometry

\end{document}

% Local Variables:
%   compile-command:
%    "rubber -d --warn all resume.tex"
% End:
